\phase{4 --- Reference CPU provider}

\textbf{Objective.} Implement and verify a CPU implementation of the generic execution contract before introducing operation algorithms. This proves the provider lifecycle independently from image-operation semantics.

\deliverables
\begin{enumerate}
  \item CPU implementation of \GEP{} with resource allocation, host access, submission, completion, synchronization, cancellation, and deterministic release.
  \item A minimal test-only work submission path sufficient to exercise the generic provider contract without introducing image-specific types into \GEP{}.
  \item Provider conformance tests for resource ownership, concurrent access, errors, cancellation, and completion ordering.
  \item Memory and allocation instrumentation that can later be reused by image operation benchmarks.
\end{enumerate}

\exitcriteria
\begin{itemize}
  \item The CPU provider passes the complete generic \GEP{} conformance suite.
  \item Resource lifetime tests prove that submitted work cannot outlive its input allocations.
  \item Cancellation and execution failure release temporary resources and report structured errors.
  \item The CPU provider remains image-agnostic and does not contain Gaussian blur, resize, or color-correction logic.
\end{itemize}

\textbf{Gate.} Close the generic provider lifecycle before adding operation implementations. This phase proves execution infrastructure, not image-processing correctness.

\phase{5 --- Core operation suite and CPU reference implementations}

\textbf{Objective.} Freeze the first useful set of image operations, implement canonical CPU versions, and create the conformance harness used by every later backend.

\deliverables
\begin{enumerate}
  \item A semantic specification for each initial operation: pixel/color model, precision, ranges, alpha treatment, border behavior, interpolation, and output type.
  \item Canonical CPU implementations for a deliberately small first set: explicit pixel/component conversion, resize, Gaussian filtering, color correction, and compositing/blending.
  \item Reusable operation conformance suites and golden/reference fixtures, including normal, boundary, degenerate, and invalid inputs.
  \item Separate tests for raster, mask, label-map, and annotation semantics where relevant.
  \item Composition tests for sequential pipelines, output/input compatibility, and failure propagation.
  \item Reference benchmarks recording CPU latency, allocations, and memory behavior.
\end{enumerate}

\exitcriteria
\begin{itemize}
  \item Every initial operation has a fully specified contract and passes its CPU reference tests.
  \item Numerical comparisons use declared absolute/relative tolerances or domain-specific metrics.
  \item Label maps are never processed with intensity interpolation unless an explicit label-aware algorithm is specified.
  \item Color calculations specify working space, transfer function, and alpha convention.
  \item Source pixels remain unchanged for non-mutating application; layer targeting preserves unrelated layers.
  \item Test fixtures are reusable without alteration by every later backend.
\end{itemize}

\textbf{Gate.} No accelerated implementation is considered correct until it passes these same fixtures and operation contracts.

\phase{6 --- ImageProcessingRuntime: qualification, selection, and execution}

\textbf{Objective.} Integrate the two protocol libraries and the proven CPU operation set into one execution path without moving image-domain qualification into the generic provider library.

\deliverables
\begin{enumerate}
  \item Registry of operation implementations and associated generic provider requirements.
  \item Runtime enumeration of available providers, devices, execution contexts, and registered implementations.
  \item Invocation of the \IPP{} qualifying function using operation instance, input/layer descriptor, implementation declaration, execution capability facts, and request constraints.
  \item Structured candidate list with accepted/rejected reasons and required preparation steps.
  \item Selection policy that only considers qualified implementations; policy inputs include device affinity, precision constraints, residency, memory budget, and caller preferences.
  \item Initial execution planner that can select one implementation for a simple operation and a compatible implementation sequence for a composed pipeline.
  \item Atomic image replacement transaction: compute and validate output first, then publish a new image revision.
  \item Observable diagnostics for the chosen implementation, provider/device, conversions, synchronization, and fallback decisions.
\end{enumerate}

\exitcriteria
\begin{itemize}
  \item An unsupported operation is rejected before execution with a useful reason.
  \item A valid implementation is selected when several candidates are registered.
  \item A faster but semantically invalid implementation is never selected.
  \item If all implementations are rejected, the runtime returns a structured failure and does not silently downgrade semantics.
  \item A failed execution leaves the input image revision unchanged.
  \item Tests prove that qualification and selection are separate and independently testable.
  \item The CPU implementations from Phase 5 are discoverable and executable through the same runtime path used by future providers.
\end{itemize}

\textbf{Gate.} Do not add heuristic routing based only on operating system or GPU vendor. Qualification must use actual operation, input, implementation, and device facts.
